\documentclass[10pt]{article}
\usepackage{fullpage}
\usepackage{graphicx}
\begin{document}
	\begin{flushright}
	Lindsey Bieda and Joe Frambach\\
	Greedy Algorithm Problems\\
	9.8.2011
	\end{flushright}
	\noindent
	6.  Consider the following problem. The input is a collection $A = \{a_{1}, \ldots , a_{n} \}$ of n points on the real line.
	The problem is to find a minimum cardinality collection $S$ of unit intervals that cover every point in
	$A$. Another way to think about this same problem is the following. You know a collection of times ($A$)
	that trains will arrive at a station. When a train arrives there must be someone manning the station.
	Due to union rules, each employee can work at most one hour at the station. The problem is to find a
	scheduling of employees that covers all the times in $A$ and uses the fewest number of employees.
	\begin{enumerate}
		\item[(a)] Prove or disprove that the following algorithm correctly solves this problem. Let $I$ be the interval
		that covers the most number of points in $A$. Add $I$ to the solution set $S$. Then recursively continue on the points in 
		$A$ not covered by $I$.
		% answer here
		\\
		\\
		Counter-example: 
		
		\begin{figure}[h]
			\centering
				\caption{Solid Line is algorithm 6a and Dotted represents optimal}
				\includegraphics{counterexample7a.png}
		\end{figure}

		\item[(b)] Prove or disprove that the following algorithm correctly solves this problem. Let $a_{j}$ be the smallest
		(leftmost) point in $A$.  Add the interval $I = (a_{j} , a_{j} + 1)$ to the solution set $S$.  Then recursively
		continue on the points in $A$ not covered by $I$.
		%answer here
		\\
		Prove that algorithm 6b is correct.
		\\
		\\
		Proof: $\exists$ input I $\ni$ algorithm 6b will produce incorrect output.\\
		\\
		$GRE(I) = [(g_{0},g_{0} + i), \ldots, (g_{n}, g_{n} + i)]$, where $i$ is the interval length\\
		$OPT(I) = [(o_{0},o_{0} + i), \ldots, (o_{n}, o_{n} + i)]$, where OPT(I)'s output is the one that agrees with GRE(I) for the most steps, $k$\\
		Let the first point of disagreement be $(o_{k}, o_{k}+1)$, therefore because greedy selected $g_{k}$ we know that $o_{k} < g_{k}$, otherwise 
		OPT would not cover point $g_{k}$ and would not be a correct solution.\\ 
		$OPT^{\prime} = OPT(I) - (o_{k}, o_{k} + i) + (g_{k}, g_{k} + i)$, $OPT^{\prime}$ retains optimal cardinality \\
		The space between $o_{k}$ and $g_{k}$ contains no points, since optimal agreed with greedy up to $k$.\\
		The inclusion space between $o_{k}+i$ and $g_{k}+i$ may only increase $OPT^{\prime}$'s cover, therefore $OPT \leq OPT^{\prime}$.\\
		Therefore, $OPT \leq OPT^{\prime} \leq OPT^{\prime\prime} \leq \ldots = GRE \bot$ 
	\end{enumerate}
	7. Consider the following problem.  The input consists of the lengths $\ell_{1}, \ldots , \ell_{n}$, and access probabilities
		$p_{1} , \ldots , p_{n}$, for $n$ files $F_{1} , \ldots , F_{n}$.  The problem is to order these files on a tape so as to minimize the
		expected access time. If the files are placed in the order $F_{s(1)} , \ldots , F_{s(n)}$ then the expected access time
		is
		
	\[E(time) = \sum\limits_{i=1}^n P_{s(i)} \sum\limits_{j=1}^{s(i)} \ell_{s(j)}\]
	\noindent
	For each of the below algorithms, either give a proof that the algorithm is correct, or a proof that the
	algorithm is incorrect.
	\begin{enumerate}
		\item[(a)] Order the files from shortest to longest on the tape. That is, $\ell_{i} < \ell_{j}$ implies that $s(i) < s(j)$.
		%answer here
		
	\[ I = \left[
		\begin{array}{ccc}
		p_{0} & \ldots & p_{n}\\ 
		\ell_{0} & \ldots & \ell_{n}
		\end{array}
	 \right]\]
	 
	 Counter-example: 
	 
	\[ I = \left[
		\begin{array}{cc}
		0.01 &  0.99\\ 
		4 & 5 
		\end{array}
	 \right]\]
	 
	 Algorithm 7a = 8.95, however, optimal produces 5.04
	 
		\item[(b)] Order the files from most likely to be accessed to least likely to be accessed.  That is, $p_{i} < p_{j}$
		implies that $s(i) > s(j)$.
		%answer here
	 Counter-example: 
	 
	\[ I = \left[
		\begin{array}{cc}
		0.5001 &  0.4999\\ 
		1000 & 1 
		\end{array}
	 \right]\]
	 	
	 	Algorithm 7a = 1000.4999, however, optimal produces 501.1
	 		
		\item[(c)] Order the the files from smallest ratio of length over access probability to largest ratio of length
		over access probability. That is, $\frac{\ell_{i}}{p_{i}} < \frac{\ell_{j}}{p_{j}}$ implies that $s(i) < s(j)$.
		%answer here
		\\
		\\
		Prove that algorithm 7c is correct. 
		\\
		Proof: $\exists$ input I $\ni$ algorithm 7c will produce incorrect output.\\
		$GRE(I) = \left[
		\begin{array}{ccc}
		gp_{0} & \ldots & gp_{n}\\ 
		g\ell_{0} & \ldots & g\ell_{n}
		\end{array}
	 \right]$ 
	 \\
	 \\
		$OPT(I) = \left[
		\begin{array}{ccc}
		op_{0} & \ldots & op_{n}\\ 
		o\ell_{0} & \ldots & o\ell_{n}
		\end{array}
	 \right]$ where OPT(I)'s output is the one that agrees with GRE(I) for the most steps, $k$\\
	 \\
	 Let the first point of disagreement be point $k$, where $\frac{\ell_{k}}{p_{k}} > \frac{\ell_{k+1}}{p_{k+1}}$, therefore $\ell_{k}~p_{k+1} > \ell_{k+1}~p_{k}$\\
	 
	 $E(time) = p_{0}\ell_{0} + \ldots + p_{k}(\ell_{0}+\ldots+\ell_{k}) + \underbrace{p_{k+1}(\ell_{0}+\ldots+\ell_{k}+\ell_{k+1})}{} + \ldots + p_{n}(\ell_{0}+\ldots+\ell_{n})$\\
	 \\
	 $p_{k+1}(\ell_{0}+\ldots+\ell_{k}+\ell_{k+1})$\\
	 $= p_{k+1}(\ell_{0}+\ldots+\ell_{k-1}) + p_{k+1}~\ell_{k}+p_{k+1}~\ell_{k+1}$\\
	 $> p_{k+1}(\ell_{0}+\ldots+\ell_{k-1}) + p_{k}~\ell_{k+1}+p_{k+1}~\ell_{k+1}$
	 \\
	 This shows that any modification made to OPT results in a larger average access time. Therefore, OPT is not optimal. $\bot$
	 
	\end{enumerate}
\end{document}